\section{Hipotesis dan Premis}

Penelitian ini dibangun di atas premis bahwa penerapan ABAC berbasis \textit{edge} pada lingkungan IoT tidak hanya bergantung pada kemampuan model dalam menyediakan kontrol akses yang \textit{fine-grained} dan \textit{context-aware}, tetapi juga pada cara atribut dikelola selama proses evaluasi akses. Pengambilan atribut secara langsung dari PIP atau \textit{cloud} dapat menjaga kemutakhiran atribut, tetapi meningkatkan latensi dan jumlah \textit{PIP calls}. Sebaliknya, \textit{cache} atribut pada \textit{gateway edge} dapat menurunkan latensi, tetapi berisiko menghasilkan \textit{stale atribute} yang dapat menyebabkan \textit{false allow, false deny,} atau keterlambatan pencabutan akses.

Berdasarkan premis tersebut, penelitian ini mengusulkan \textit{GWO-optimized ABAC}, yaitu mekanisme ABAC berbasis \textit{edge} yang menggunakan GWO untuk menentukan konfigurasi \textit{cacheable attribute,} TTL per atribut, dan ukuran \textit{cache} lokal pada \textit{gateway}. Hipotesis penelitian dirumuskan sebagai berikut:

\begin{description}
    \item [H1 -- Efisiensi performa \textit{GWO-optimized ABAC}.] \textit{GWO-optimized ABAC} mampu menurunkan latensi keputusan dan jumlah \textit{PIP calls} dibandingkan \textit{static cache ABAC, adaptive TTL ABAC,} dan \textit{random search}.
    \item [H2 -- Kualitas keputusan otorisasi \textit{GWO-optimized ABAC}.] \textit{GWO-optimized ABAC} mampu mempertahankan kualitas keputusan otorisasi dengan menjaga \textit{false allow rate, false deny rate, stale cache hit rate,} dan \textit{revocation false allow rate} pada tingkat yang lebih terkendali dibandingkan pendekatan \textit{cache} berbasis konfigurasi statis atau manual.
    \item [H3 -- Keunggulan konfigurasi hasil GWO.] Konfigurasi \textit{cacheable attribute}, TTL per atribut, dan ukuran \textit{cache} yang dihasilkan oleh GWO lebih efektif dibandingkan konfigurasi manual dan \textit{random search} dalam menyeimbangkan efisiensi \textit{cache, attribute freshness}, dan risiko keputusan akses.
    \item [H4 -- Ketahanan \textit{GWO-optimized ABAC} terhadap \textit{workload} dinamis.] \textit{GWO-optimized ABAC} mampu mempertahankan performa dan kualitas keputusan akses pada skenario \textit{worklaod} rendah, sedang, tinggi, dan sangat tinggi yang melibatkan perubahan \textit{trust, anomaly,} kondisi jaringan, status \textit{firmware, role, critical request, PIP unavailability,} dan \textit{revocation event}.
    \item [H5 -- \textit{Trade-off} performa dan keamanan pada \textit{GWO-optimized ABAC}.] \textit{GWO-optimized ABAC} mampu menghasilkan \textit{trade-off} yang lebih seimbang antara penurunan latensi dan pengendalian risiko \textit{stale attribute} dibandingkan pendekatan yang hanya menekankan \textit{cache hit} atau TTL statis.
\end{description}

Secara formal, intuisi dasar penelitian ini dapat dinyatakan sebagai berikut. Jika keuntungan performa dari penggunaan \textit{cache} atribut dinyatakan sebagai $G$, sedangkan risiko keamanan akibat \textit{stale attribute} dinyatakan sebagai $R_s$, maka konfigurasi \textit{cache} layak digunakan ketika $G > R_s$. Namun karena \textit{false allow} merupakan pelanggaran keamanan yang lebih kritis daripada sekadar peningkatan latensi, tujuan \textit{GWO-optimized ABAC} bukan hanya memaksimalkan \textit{cache hit} atau menurunkan \textit{PIP calls}, tetapi mencari konfigurasi yang mampu mendorong sistem dari kondisi latensi rendah tetapi risiko \textit{stale} tinggi menuju kondisi latensi rendah dengan risiko keputusan akses terkendali. Dengan demikian, target utama penelitian ini adalah \textit{GWO-optimized ABAC} mampu memberikan keseimbangan yang lebih baik antara efisiensi evaluasi akses, \textit{attribute freshness}, dan keamanan keputusan otorisasi dibandingkan pendekatan ABAC berbasis \textit{edge} dengan konfigurasi manual atau \textit{non-optimized}.